\section{Identification by Proximal Balancing}
\label{sec:identification}

We introduce \emph{proximal balancing}, an identification method that brings covariate-balancing ideas \citep{hainmueller2012entropy,imai2014covariate} to high-dimensional pretreatment proxies. It learns a representation from the retained proxy blocks and uses the held-out block with observed covariates to audit residual treatment information.

\subsection{Exact Balance}
\label{sec:probe}

We first consider a fixed held-out block within the high-dimensional proxy $W=(W_1,\ldots,W_J)$. For $j\in[J]$, Definition~\ref{def:held-out-remaining-blocks} calls $W_j$ the \emph{held-out block} and $W_{-j}$ the \emph{remaining blocks}.

We begin with two structural conditions: $(U,V_j)=(U,X,W_{-j})$ suffices to remove treatment confounding for each potential outcome, and $W_j$ is independent of treatment given $(U,V_j)$.
\begin{assumption}[Latent exchangeability and held-out treatment independence]
\label{ass:joint-latent-exchangeability}\label{ass:held-out-separation}
	For each $a\in\{0,1\}$, $Y(a)\perp\!\!\!\perp A\mid(U,V_j)$ and $W_j\perp\!\!\!\perp A\mid(U,V_j)$.
\end{assumption}
Latent exchangeability says that, after $(U,V_j)$ is fixed, treatment carries no further information about $Y(a)$. Held-out treatment independence states that $W_j$ carries no additional treatment information given the latent and retained variables.

Let $\phi_j$ be a measurable encoder and set $Z_j\triangleq\phi_j(V_j)=\phi_j(X,W_{-j})$. Proximal balancing seeks a representation that makes the treatment arms comparable for the potential-outcome means, while $T_j=(X,W_j)$ supplies an observable check of residual treatment information.  For a distribution $Q$ of $(V_j,U)$, let $\mathcal K_jQ$ be the distribution of $T_j$ obtained by drawing $(V_j,U)\sim Q$ and then $W_j\sim P(W_j\mid V_j,U)$. Let $\mu_{a,j,z}(v,u)\triangleq\mathbb E\{Y(a)\mid V_j=v,U=u,Z_j=z\}$ be the mean of $Y(a)$ at a latent state.




Different  distributions $Q,Q'$ of $(V_j,U)$ can induce the same  distribution of $T_j$. The next assumption states when this loss of latent information is harmless for the potential-outcome means needed for identification.

\begin{assumption}[Outcome-relevant completeness]
\label{ass:outcome-relevant-completeness}
For  almost every $z$, every  pair of distributions $Q,Q'$ of $(V_j,U)$ that are absolutely continuous with respect to $P(V_j,U\mid Z_j=z)$ and under which $\mu_{0,j,z}$ and $\mu_{1,j,z}$ are integrable, and each $a\in\{0,1\}$, $\mathcal K_jQ=\mathcal K_jQ'$ implies $\mathbb E_Q\mu_{a,j,z}=\mathbb E_{Q'}\mu_{a,j,z}$.
\end{assumption}

 In words, if two distributions of the latent state $Q,Q'$ give different means of $Y(a)$, they also give different distributions of $T_j=(X,W_j)$.

The following theorem combines these structural conditions with overlap and exact held-out balance to identify the average treatment effect.

\begin{theorem}[Identification from a held-out audit]
\label{thm:exact-proxy-identification}
Fix $j\in[J]$ and a measurable encoder $\phi_j$, and set $Z_j\triangleq\phi_j(V_j)$. Suppose Assumptions~\ref{ass:consistency}--\ref{ass:outcome-relevant-completeness} hold. If
\begin{align}
	0 &< P(A=1\mid Z_j)<1\ \text{a.s.},\qquad (X,W_j)\perp\!\!\!\perp A\mid Z_j
\label{eq:exact-proxy-balance}
\end{align}
then
\begin{align}
\tau &=\mathbb E[\mathbb E(Y\mid A=1,Z_j)-\mathbb E(Y\mid A=0,Z_j)].
\label{eq:exact-identification}
\end{align}
\end{theorem}

Under Assumptions~\ref{ass:consistency}--\ref{ass:outcome-relevant-completeness}, Thm.~\ref{thm:exact-proxy-identification} shows that the average treatment effect is identified by Eq.~\eqref{eq:exact-identification} whenever $Z_j=\phi_j(V_j)$ satisfies the overlap and exact held-out balance conditions in Eq.~\eqref{eq:exact-proxy-balance}.

\subsection{Unknown Valid Block}
\label{sec:unknown-blocks}

Section~\ref{sec:probe} treats the ideal case in which a valid held-out block is known and exact balance holds. We next allow the valid block to be unknown and then quantify the bias caused by imperfect balance.

Let $\mathfrak S\triangleq\{S\subset[J]:1\le |S|\le J-1\}$ be the collection of nonempty proper subsets of proxy-block indices. For $S\in\mathfrak S$, write $W_S=(W_k)_{k\in S}$ for the held-out blocks and $W_{-S}=(W_k)_{k\notin S}$ for the remaining blocks, set $V_S\triangleq(X,W_{-S})$, and let $Z_S\triangleq\phi_S(V_S)$. The singleton $S=\{j\}$ recovers the split of Section~\ref{sec:probe}. Each split yields the population candidate adjustment value
\begin{align}
\label{eq:block-candidate-adjustment-value}
	\theta_S &\triangleq\mathbb E[\mathbb E(Y\mid A=1,Z_S)-\mathbb E(Y\mid A=0,Z_S)].
\end{align}

The next definition records the observable screen and the target plurality condition.
\begin{definition}[Proximal plurality]
\label{def:population-plurality}
Let $\mathcal B$ be the collection of splits $S\in\mathfrak S$ for which overlap $0<P(A=1\mid Z_S)<1$ holds almost surely and exact balance $(X,W_S)\perp\!\!\!\perp A\mid Z_S$ holds. The target has a unique proximal plurality when
\begin{align}
|\{S\in\mathcal B:\theta_S=\tau\}| &>\max_{c\ne\tau:\{S\in\mathcal B:\theta_S=c\}\ne\varnothing}|\{S\in\mathcal B:\theta_S=c\}|,
\label{eq:unique-target-plurality}
\end{align}
where the maximum over an empty set is zero.
\end{definition}
In words, proximal plurality means that $\tau$ is the most frequent candidate value among the screened splits, even if it is not attained by a majority. For example, let $J=3$, so $\mathfrak S=\{\{1\},\{2\},\{3\},\{1,2\},\{1,3\},\{2,3\}\}$. If $\mathcal B=\{\{1\},\{2\},\{1,2\},\{1,3\}\}$, $\tau=1$, $\theta_{\{1\}}=\theta_{\{1,2\}}=1$, $\theta_{\{2\}}=0.5$, and $\theta_{\{1,3\}}=0.3$, then $\tau$ appears twice while each competing value appears once, so the target has a unique proximal plurality.

When $\mathcal B\ne\varnothing$ and proximal plurality holds, the causal effect is identified by the population mode of the candidate values.
\begin{theorem}[Identification with an unknown valid block]
\label{thm:unknown-valid-block-identification}
If $\mathcal B\ne\varnothing$ and the target has a unique proximal plurality, then $\{\tau\}=\operatorname*{arg\,max}_{c\in\mathbb R}|\{S\in\mathcal B:\theta_S=c\}|$.
\end{theorem}

Invalid passing blocks may remain. Identification requires the correct candidates to form the unique largest equality class in Eq.~\eqref{eq:unique-target-plurality}.

\subsection{Approximate Balance}
\label{sec:imperfect-balance}

Exact balance in Eq.~\eqref{eq:exact-proxy-balance} may be infeasible in practice. We derive the bound for a fixed held-out block $W_j$. The same argument extends to a grouped held-out block $W_S$ after replacing the single-block objects by their grouped counterparts. Define $Z_j^\phi \triangleq \phi(V_j)$, $e_{j,\phi}(z) \triangleq P(A=1\mid Z_j^\phi=z)$ and $q_{j,\phi}(t,z) \triangleq P(A=1\mid T_j=t,Z_j^\phi=z)$.

\begin{definition}[Residual treatment discrepancy]
\label{def:held-out-proxy-discrepancy}
The residual treatment discrepancy of $\phi$ is
\begin{align}\label{eq:residual-treatment-discrepancy}
	D_{\mathrm{res},j}(\phi) &\triangleq\left[\mathbb E\bigl[\{q_{j,\phi}(T_j,Z_j^\phi)-e_{j,\phi}(Z_j^\phi)\}^2\bigr]\right]^{1/2}.
\end{align}
\end{definition}
The discrepancy $D_{\mathrm{res},j}(\phi)$ is the $L_2$ magnitude of the additional treatment information in $T_j=(X,W_j)$ beyond $Z_j^\phi$. It characterizes exact conditional balance as follows.

\begin{proposition}[Zero discrepancy]
\label{prop:zero-discrepancy}
	$D_{\mathrm{res},j}(\phi)=0$ if and only if $A\perp\!\!\!\perp(X,W_j)\mid Z_j^\phi$.
\end{proposition}

To translate residual treatment discrepancy into bias of an adjusted effect, we impose the following stability condition.
\begin{assumption}[$L_2$ outcome-relevant stability]
\label{ass:uniform-stability}
For each encoder $\phi$, let $\mu^\phi_{a,j,z}$ be $\mu_{a,j,z}$ with $Z_j^\phi = z$, let $P_{\phi,z}\triangleq P(T_j\mid Z_j^\phi=z)$, and for a signed measure $\nu$ with a density with respect to $P_{\phi,z}$ write $\|\nu\|_{\phi,z}\triangleq\|d\nu/dP_{\phi,z}\|_{L_2(P_{\phi,z})}$. There is a finite $\Gamma_j(\phi)$ such that, for almost every $z$, each $a\in\{0,1\}$, and all $Q,Q'$ as in Assumption~\ref{ass:outcome-relevant-completeness} with $Z_j^\phi$ in place of $Z_j$,
\[|\mathbb E_Q\mu^\phi_{a,j,z}-\mathbb E_{Q'}\mu^\phi_{a,j,z}|\le\Gamma_j(\phi)\,\|\mathcal K_jQ-\mathcal K_jQ'\|_{\phi,z}.\]
\end{assumption}
 If $\mathcal K_jQ=\mathcal K_jQ'$, the right side of Assumption~\ref{ass:uniform-stability} is zero, so Assumption~\ref{ass:uniform-stability} implies Assumption~\ref{ass:outcome-relevant-completeness} for $Z_j^\phi$. It gives additional rate information: when the two distributions of $T_j$ are close, the two means of $Y(a)$ are close, within the unobserved factor $\Gamma_j(\phi)$.



Under Assumption~\ref{ass:uniform-stability}, imperfect balance yields a quantitative bound on causal bias.

\begin{theorem}[Approximate-balance bias]
\label{thm:approximate-proxy-bias-bound}
For any representation $Z$, define its ordinary adjustment value as $\tau_Z\triangleq\mathbb E[\mathbb E(Y\mid A=1,Z)-\mathbb E(Y\mid A=0,Z)]$.
Under Assumptions~\ref{ass:consistency}, \ref{ass:joint-latent-exchangeability}, and \ref{ass:uniform-stability}, and overlap $\eta\le P(A=1\mid Z_j^\phi)\le1-\eta$ almost surely for some $\eta\in(0,1/2]$,
\begin{align}
|\tau_{Z_j^\phi}-\tau| &\le\frac{\Gamma_j(\phi)}{\eta(1-\eta)}D_{\mathrm{res},j}(\phi).
\label{eq:uniform-approximate-proxy-bias-bound}
\end{align}
\end{theorem}

Eq.~\eqref{eq:uniform-approximate-proxy-bias-bound} separates observable discrepancy from unobserved sensitivity. A small discrepancy cannot verify assumptions, and $\Gamma_j(\phi)$ requires substantive analysis. Section~\ref{sec:estimation} adds estimation errors.
